\section{Conclusions and outlook}
Stationary voltage inversion reconstructs flux without enforcing magnetic
symmetry. Under the stated reflection-symmetric quasi-static assumptions,
d flux is even and q flux odd in quadrature current, including saturation and
cross-saturation. Their forbidden half-difference and half-sum are useful
observables because they preserve information that a symmetric average
would discard. They diagnose incompatibility with the baseline, rather than
uniquely identifying a fault.

The resistance-equivalent quantity $\resistanceequivalent=-\omegae r_d/a_q$
is explicitly conditional: it equals used-minus-true resistance only for a
sole resistance mismatch with adequate matching. In the outlier experiment,
all three 2000-rpm slices exhibit substantial operating-point dependence.
The high-current subset is empirically flatter, with medians $0.096$,
$-0.579$ and $-1.268$ m$\Omega$ at rotor references 40, 60 and 80 $^\circ$C.
This is a thermal-reference comparison, not quantitative winding thermometry.
In the separate multi-speed experiment, six and seven high-current pairs
yield $-2.157$ and $-5.334$ m$\Omega$ at 1000 and 3000 rpm. At 36 common
valid keys the direct inverse-speed prediction has an RMS discrepancy of
$0.1573$ mWb, essentially unchanged by accounting for measured pair-current
magnitudes. One constant resistance mismatch is insufficient to explain the
full observed structure under the specified baseline. A resistance-like
constant may describe a subset, but its causal share is not identified.

The exact angle transformation rotates both current arguments and flux
vectors. It preserves reciprocity while rotating the natural symmetry axis;
saturation changes its sensitivity through the differential inductance
matrix. A current-proportional voltage error can exactly imitate resistance
in reconstruction. The retained synthetic checks verify these derivations
and correction signs, not measured parameter recovery. Unique physical
separation into resistance, iron-loss, angle and inverter contributions
remains open and requires independent thermal, voltage, current and magnetic
information.

The conservation cross-check strengthens this conclusion. The symmetry and
path equivalents differ at 2000 rpm/60 $^\circ$C even after exact P1
integration. A symmetry-based resistance hypothesis improves the selected
high-current region but worsens all-pair d/vector residuals and relative
path disagreement. Its invariant path equivalent is an algebraic check,
not an independent resistance calibration. Interpolation and terminal-current
model reduction remain limits; a smaller selected symmetry residual cannot
justify automatic correction.

\paragraph{Outlook.}
A future joint identification should combine common-support speed sweeps,
independent winding-temperature measurements, and additional calibrated
voltage, power, or magnetic-state observables. Speed diversity tests a sole
resistance mismatch but cannot resolve a voltage error with the same
inverse-speed fingerprint. Storage, loss, and measurement terms need distinct
observation relations and a demonstrable parameter rank.
The complementary question is how flux should be represented for the
conservative magnetic storage subsystem in suitable conjugate coordinates, with
$\flux=\nabla_{\current}\coenergy$ and
$\Ldiff=\nabla^2_{\current}\coenergy$. That question belongs to the separate
co-energy reconstruction draft \cite{companion-coenergy}. Fitting inconsistent
terminal maps would be a model projection with a retained diagnostic residual;
no such real-data fit or joint causal identification is carried out here.
